\subsection{RQ1: End-to-End Effectiveness}
\label{sec:rq1}

% Notes: docs/exhibits.md#tab_rq1_end_to_end_summary
\begin{table*}[t]
  \centering
  \setlength{\abovecaptionskip}{2pt}%
  \setlength{\belowcaptionskip}{3pt}%
  \caption{RQ1 baseline failure reference and our-controller behavior.}
  \label{tab:rq1_preventive_alignment}
  \label{tab:rq1_controller_behavior}

  {\scriptsize
  \setlength{\arrayrulewidth}{0.35pt}%
  \setlength{\doublerulesep}{0.8pt}%
  \renewcommand{\arraystretch}{1.08}%
  \newlength{\rqsurvivedwidth}%
  \settowidth{\rqsurvivedwidth}{\textbf{10/10} (--)}%
  \newcommand{\rqsurvived}[1]{\makebox[\rqsurvivedwidth][r]{#1}}%
  \newsavebox{\rqtabbox}%
  \newcommand{\rqtabular}{%
  \begin{tabular}{|c|c|c|r||c||c||rr|rr||rr|rr|}
    \hline\hline
    \multirow{4}{*}[-1.8ex]{\textbf{Device}} &
    \multirow{4}{*}[-1.8ex]{\textbf{Condition}} &
    \multirow{4}{*}[-1.8ex]{\makecell[c]{\textbf{Starting}\\\textbf{overheat}\\\textbf{level}}} &
    \multirow{4}{*}[-1.8ex]{\textbf{N}} &
    \multicolumn{1}{c||}{\multirow{2}{*}{\makecell[c]{\textbf{Baseline}}}} &
    \multicolumn{9}{c|}{\textbf{Our-controller behavior}} \\
    \cline{6-14}
    & & & &
    \multicolumn{1}{c||}{} &
    \multicolumn{1}{c||}{\textbf{Deadline safety}} &
    \multicolumn{4}{c||}{\textbf{Draft stages executed}} &
    \multicolumn{4}{c|}{\textbf{Pacing cost}} \\
    \cline{5-5}\cline{6-14}
    & & & &
    \multicolumn{1}{c||}{\multirow{2}{*}[-0.4ex]{%
      \makecell[c]{\textbf{Survived runs}\\\textbf{(timeout onset)}}}} &
    \multicolumn{1}{c||}{\multirow{2}{*}[-0.4ex]{%
      \makecell[c]{\textbf{Survived runs}\\\textbf{(timeout onset)}}}} &
    \multicolumn{2}{c|}{\makecell[c]{\mbox{\boldmath$M$}\\\textbf{(captures)}}} &
    \multicolumn{2}{c||}{\makecell[c]{\mbox{\boldmath$S$}\\\textbf{(captures)}}} &
    \multicolumn{2}{c|}{\makecell[c]{\textbf{Activated}\\\textbf{(transitions)}}} &
    \multicolumn{2}{c|}{\makecell[c]{\mbox{\boldmath$d$}\textbf{ P50}\\\textbf{(ms)}}} \\
    \cline{7-14}
    & & & & & &
    \multicolumn{1}{c}{\textbf{@5}} & \multicolumn{1}{c|}{@30} &
    \multicolumn{1}{c}{\textbf{@5}} & \multicolumn{1}{c||}{@30} &
    \multicolumn{1}{c}{\textbf{@5}} & \multicolumn{1}{c|}{@30} &
    \multicolumn{1}{c}{\textbf{@5}} & \multicolumn{1}{c|}{@30} \\
    \hline\hline
    \multirow{14}{*}{\makecell[c]{S26\\Ultra}}
      & \multirow{7}{*}{\makecell[c]{12MP\\normal}}
      & Lv0 & 10 & \rqsurvived{\textbf{10/10} (--)} & \rqsurvived{\textbf{10/10} (--)}
        & \textbf{5.0} & 30.0 & \textbf{5.0} & 30.0
        & \textbf{0.0} & 0.0 & \textbf{--} & -- \\
    \cline{3-14}
      & & Lv1 & 10 & \rqsurvived{\textbf{10/10} (--)} & \rqsurvived{\textbf{10/10} (--)}
        & \textbf{5.0} & 29.4 & \textbf{5.0} & 30.0
        & \textbf{0.0} & 0.2 & \textbf{--} & 127 \\
    \cline{3-14}
      & & Lv2 & 10 & \rqsurvived{7/10 (24)} & \rqsurvived{\textbf{10/10} (--)}
        & \textbf{5.0} & 27.1 & \textbf{5.0} & 30.0
        & \textbf{0.0} & 3.6 & \textbf{--} & 262 \\
    \cline{3-14}
      & & Lv3 & 10 & \rqsurvived{1/10 (13)} & \rqsurvived{\textbf{10/10} (--)}
        & \textbf{5.0} & 25.6 & \textbf{5.0} & 29.4
        & \textbf{0.0} & 10.4 & \textbf{--} & 284 \\
    \cline{3-14}
      & & Lv4 & 10 & \rqsurvived{0/10 (8)} & \rqsurvived{\textbf{10/10} (--)}
        & \textbf{5.0} & 17.4 & \textbf{5.0} & 28.7
        & \textbf{0.0} & 11.0 & \textbf{--} & 437 \\
    \cline{3-14}
      & & Lv5 & 10 & \rqsurvived{0/10 (8)} & \rqsurvived{\textbf{10/10} (--)}
        & \textbf{5.0} & 13.6 & \textbf{5.0} & 28.3
        & \textbf{0.6} & 12.2 & \textbf{547} & 470 \\
    \cline{3-14}
      & & Lv6 & 10 & \rqsurvived{0/10 (7)} & \rqsurvived{\textbf{10/10} (--)}
        & \textbf{4.7} & 10.9 & \textbf{4.8} & 26.5
        & \textbf{0.9} & 9.9 & \textbf{503} & 521 \\
    \Xcline{2-14}{0.7pt}
      & \multirow{7}{*}{\makecell[c]{24MP\\memory\\pressure}}
      & Lv0 & 10 & \rqsurvived{7/10 (26)} & \rqsurvived{\textbf{10/10} (--)}
        & \textbf{5.0} & 27.8 & \textbf{5.0} & 30.0
        & \textbf{0.2} & 11.0 & \textbf{223} & 338 \\
    \cline{3-14}
      & & Lv1 & 10 & \rqsurvived{2/10 (19)} & \rqsurvived{\textbf{10/10} (--)}
        & \textbf{5.0} & 20.0 & \textbf{5.0} & 27.0
        & \textbf{0.0} & 10.2 & \textbf{--} & 470 \\
    \cline{3-14}
      & & Lv2 & 10 & \rqsurvived{0/10 (14)} & \rqsurvived{\textbf{10/10} (--)}
        & \textbf{5.0} & 17.5 & \textbf{5.0} & 23.7
        & \textbf{0.0} & 10.8 & \textbf{--} & 377 \\
    \cline{3-14}
      & & Lv3 & 10 & \rqsurvived{0/10 (5)} & \rqsurvived{\textbf{10/10} (--)}
        & \textbf{4.7} & 10.5 & \textbf{4.7} & 16.3
        & \textbf{1.0} & 11.0 & \textbf{616} & 418 \\
    \cline{3-14}
      & & Lv4 & 10 & \rqsurvived{0/10 (3)} & \rqsurvived{\textbf{10/10} (--)}
        & \textbf{4.7} & 13.4 & \textbf{5.0} & 23.1
        & \textbf{1.5} & 10.6 & \textbf{458} & 400 \\
    \cline{3-14}
      & & Lv5 & 10 & \rqsurvived{0/10 (4)} & \rqsurvived{\textbf{10/10} (--)}
        & \textbf{4.8} & 5.2 & \textbf{5.0} & 15.8
        & \textbf{1.6} & 8.5 & \textbf{737} & 367 \\
    \cline{3-14}
      & & Lv6 & 10 & \rqsurvived{0/10 (3)} & \rqsurvived{\textbf{10/10} (--)}
        & \textbf{4.5} & 7.5 & \textbf{4.6} & 14.0
        & \textbf{1.8} & 8.5 & \textbf{685} & 525 \\
    \Xcline{1-14}{0.7pt}
    \multirow{14}{*}{S26}
      & \multirow{7}{*}{\makecell[c]{12MP\\normal}}
      & Lv0 & 10 & \rqsurvived{7/10 (28)} & \rqsurvived{\textbf{10/10} (--)}
        & \textbf{5.0} & 30.0 & \textbf{5.0} & 30.0
        & \textbf{0.0} & 2.6 & \textbf{--} & 159 \\
    \cline{3-14}
      & & Lv1 & 10 & \rqsurvived{5/10 (26)} & \rqsurvived{\textbf{10/10} (--)}
        & \textbf{5.0} & 30.0 & \textbf{5.0} & 30.0
        & \textbf{0.0} & 3.7 & \textbf{--} & 241 \\
    \cline{3-14}
      & & Lv2 & 10 & \rqsurvived{4/10 (22)} & \rqsurvived{\textbf{10/10} (--)}
        & \textbf{5.0} & 30.0 & \textbf{5.0} & 30.0
        & \textbf{0.0} & 5.3 & \textbf{--} & 256 \\
    \cline{3-14}
      & & Lv3 & 10 & \rqsurvived{0/10 (15)} & \rqsurvived{\textbf{10/10} (--)}
        & \textbf{5.0} & 26.9 & \textbf{5.0} & 30.0
        & \textbf{0.0} & 11.8 & \textbf{--} & 277 \\
    \cline{3-14}
      & & Lv4 & 10 & \rqsurvived{0/10 (9)} & \rqsurvived{\textbf{10/10} (--)}
        & \textbf{5.0} & 16.8 & \textbf{5.0} & 22.4
        & \textbf{0.1} & 15.4 & \textbf{31} & 297 \\
    \cline{3-14}
      & & Lv5 & 10 & \rqsurvived{0/10 (9)} & \rqsurvived{\textbf{10/10} (--)}
        & \textbf{5.0} & 13.7 & \textbf{5.0} & 23.7
        & \textbf{0.2} & 12.7 & \textbf{183} & 392 \\
    \cline{3-14}
      & & Lv6 & 10 & \rqsurvived{0/10 (9)} & \rqsurvived{\textbf{10/10} (--)}
        & \textbf{5.0} & 12.4 & \textbf{5.0} & 19.9
        & \textbf{0.4} & 15.1 & \textbf{237} & 410 \\
    \Xcline{2-14}{0.7pt}
      & \multirow{7}{*}{\makecell[c]{24MP\\memory\\pressure}}
      & Lv0 & 10 & \rqsurvived{1/10 (23)} & \rqsurvived{\textbf{10/10} (--)}
        & \textbf{5.0} & 30.0 & \textbf{5.0} & 30.0
        & \textbf{0.0} & 15.2 & \textbf{--} & 330 \\
    \cline{3-14}
      & & Lv1 & 10 & \rqsurvived{1/10 (21)} & \rqsurvived{\textbf{10/10} (--)}
        & \textbf{5.0} & 28.0 & \textbf{5.0} & 28.6
        & \textbf{0.0} & 12.4 & \textbf{--} & 255 \\
    \cline{3-14}
      & & Lv2 & 10 & \rqsurvived{0/10 (19)} & \rqsurvived{\textbf{10/10} (--)}
        & \textbf{5.0} & 17.8 & \textbf{5.0} & 23.0
        & \textbf{0.0} & 14.2 & \textbf{--} & 419 \\
    \cline{3-14}
      & & Lv3 & 10 & \rqsurvived{0/10 (9)} & \rqsurvived{\textbf{10/10} (--)}
        & \textbf{5.0} & 21.5 & \textbf{5.0} & 29.1
        & \textbf{0.3} & 17.4 & \textbf{96} & 466 \\
    \cline{3-14}
      & & Lv4 & 10 & \rqsurvived{0/10 (6)} & \rqsurvived{\textbf{10/10} (--)}
        & \textbf{5.0} & 22.2 & \textbf{5.0} & 27.4
        & \textbf{0.1} & 17.7 & \textbf{31} & 458 \\
    \cline{3-14}
      & & Lv5 & 10 & \rqsurvived{0/10 (5)} & \rqsurvived{\textbf{10/10} (--)}
        & \textbf{5.0} & 7.2 & \textbf{5.0} & 9.2
        & \textbf{0.9} & 7.3 & \textbf{397} & 538 \\
    \cline{3-14}
      & & Lv6 & 10 & \rqsurvived{0/10 (6)} & \rqsurvived{\textbf{10/10} (--)}
        & \textbf{5.0} & 6.5 & \textbf{5.0} & 7.8
        & \textbf{1.5} & 6.4 & \textbf{367} & 467 \\
    \hline\hline
  \end{tabular}}%
  \fittabcolsep{\rqtabbox}{\rqtabular}{\textwidth}{28}%
  \rqtabular
  }
\end{table*}

Table~\ref{tab:rq1_controller_behavior} evaluates the full controller across 28 cells, each defined by a device, a capture condition, and a starting overheat level.
The \emph{Baseline} column provides a failure reference by executing \(M\) and \(S\) on every capture with the thermal guard bypassed.
Because Baseline runs were collected separately, the comparison is at the cell level rather than between paired runs.

\smallskip\noindent\textbf{Deadline safety.}
In all 28 cells, all runs completed the 30-capture horizon without an observed Capture Timeout.
Baseline reached Capture Timeout in 26 of the 28 cells, with the earliest failure occurring at the third capture. 
In 17 cells, all ten runs timed out.
Overall, the controller maintained timeout-free completion across all devices, conditions, and starting overheat levels over the 30-capture horizon.

\smallskip\noindent\textbf{Stage execution.}
The controller did not obtain this result by broadly suppressing optional stages.
Over the first five captures, \(M\) and \(S\) executed on at least 4.5 and 4.6 captures per run in every cell, and 23 of the 28 cells executed both stages in all five captures.
Over the full horizon, the controller executed \(M\) on 5.2--22.2 captures per run even in the twelve cells starting at overheat level 4 or higher, where the production thermal guard begins suppressing optional stages.
\(S\) was executed at least as often as \(M\) in all twelve cells.
Thus, even above the production cutoff, both optional stages remained active across substantial fractions of the 30 captures.
This pattern shows that the controller retained optional stages progressively rather than reproducing the static cutoff's all-or-nothing behavior.

\smallskip\noindent\textbf{Pacing cost.}
Pacing was infrequent during the first five captures.
It activated on none of the four eligible transitions in 14 cells and on at most 1.8 transitions per run in any cell, although \(d\)~P50 reached 737~ms when pacing did act.
Over all 30 captures, pacing activated on as many as 17.7 of the 29 eligible transitions per run, while \(d\)~P50 remained at 127--538~ms in cells with positive delays.
Thus, pacing preserved most of the responsiveness gained from parallel capture during early consecutive capture.